\section*{Résumé}
\addcontentsline{toc}{section}{Résumé}

rockim est un code de calcul par éléments finis et discrets combinés (FDEM), écrit en C++ pendant la thèse pour simuler l'impact d'un insert en carbure de tungstène sur une roche. Il appartient à la lignée de Munjiza~\cite{munjiza2004fdem}, comme Y-Geo, HOSS, MultiFracS et Solidity, et reprend en particulier le modèle de fracture 3D que Guo a développé pour Solidity~\cite{guo2014thesis}. La roche est un assemblage de tétraèdres élastiques reliés par des éléments joints cohésifs ; un joint s'endommage selon un critère de Mohr-Coulomb avec coupure en traction, puis rompt, et les fragments libérés interagissent par contact frottant.

Le rapport décrit la formulation dans l'ordre de la thèse de Guo (\cref{sec-formulation}), puis rassemble les tests à solution exacte (\cref{sec-verification}), les sensibilités numériques (\cref{sec-sensibilites}), trois reproductions de code à code, sur les essais d'insertion adaptative de Yan et al. (\cref{sec-yan}), la fracturation hydraulique d'AbuAisha et al. (\cref{sec-abuaisha}) et la zone endommagée autour d'un tunnel de Wang et al. (\cref{sec-tunnel}), les bancs élémentaires préparés (\cref{sec-validation}), l'impact d'un insert (\cref{sec-impact}), la calibration sur le granite de Red Bohus (\cref{sec-calib}), le coût (\cref{sec-cout}) et la discussion (\cref{sec-discussion}).

Les briques vérifiables isolément passent leurs tests à solution exacte : une collision par le potentiel de contact conserve l'énergie cinétique à $5\times10^{-7}$ près au pire, une onde de compression dans une barre se propage à $0{,}05~\%$ de la célérité exacte en éléments finis, un bloc glissant s'arrête à moins de $1{,}74~\%$ de la distance exacte, et le résidu du bilan d'énergie est nul à l'arrondi près en élastique. Les mêmes tests chiffrent deux biais propres à la FDEM : des joints présents dès l'origine avec une pénalité de 20 fois $E/h$ ralentissent l'onde de $2{,}9$ à $3{,}2~\%$, et le contact par pénalité nœud-face fait basculer un bloc qui devrait glisser.

Les reproductions de code à code retrouvent les grandeurs globales de leurs références. Sur les essais de Yan et al., les grandeurs comparées tombent à moins de $1~\%$ de la référence, sauf l'angle de frottement de l'insertion intrinsèque ($-3{,}8~\%$) et l'angle de rupture ($+6{,}7~\%$), mais le plateau intrinsèque reste 13 à $16~\%$ sous l'insertion adaptative. La pression de rupture du forage d'AbuAisha et al. est surestimée de $+6{,}7$ à $+25{,}0~\%$, contre $+4~\%$ pour Y-Geo, avec un faciès conforme. Autour du tunnel, le rayon de la zone endommagée de Wang et al. est retrouvé à $-4$ à $-21~\%$ entre 4 et 7~MPa, sans la texture en blocs.

Sur l'impact d'un insert sur le calcaire de St Anne, comparé à la simulation de Yang et al.~\cite{yang2025validation} et non à leur essai, quatre des sept critères tombent à 6 à $12~\%$, l'enfoncement est trop grand de $17~\%$ et la masse de fragments de $88~\%$ ($4{,}70$~g, contre $2{,}5$~g simulés et $1{,}2$~g mesurés) ; cette comparaison n'est pas indépendante du choix des paramètres. La réplique sur le granite de Kuru échoue : vitesse d'indentation trop grande de $22~\%$, pulvérisation quasi inactive. Sur le Red Bohus, aucun jeu FDEM ne restitue à la fois la résistance en compression simple et l'enveloppe confinée, et la seule prédiction hors calage, les triaxiaux à 75 et 100~MPa, échoue ($+29$ et $+36~\%$). Aucun impact n'a encore été comparé à un essai qui n'a pas servi au calage.

Trois réserves bornent ces résultats (\cref{sec-discussion}). Dans les calculs de tunnel les plus fragmentés, le travail net du contact dépasse l'énergie cohésive. Les reproductions de Yan, d'AbuAisha et du tunnel ont tourné avec des binaires d'août 2026, et la compression de Yan rejouée en octobre ne reproduit plus le résultat d'août. La réplique St Anne n'a tourné que sur un maillage, plus grossier que celui de Yang, avec des réglages qui s'écartent de la configuration de référence.
\clearpage
